\documentclass{article}

\usepackage{amsfonts}
\usepackage{amsthm}
\usepackage{mathrsfs}
\newtheorem{theorem}{Theorem}
\theoremstyle{definition}\newtheorem{definition}{DEF}
\theoremstyle{lemma}\newtheorem{lemma}{Lemma}
\theoremstyle{corrolary}\newtheorem{corrolary}{Cor}

\begin{document}

Many processes will provide different results upon repitition with the same inputs, which we call randomness.
Statistics is the study of randomness and random expirements, but interestingly there is no good definition of randomness.
Instead, we simply say a random expirement is any process that exhibits this behavior, and do not concern ourselves with the particulars of what causes that randomness.

\begin{definition}
Given $n \in \mathbb{N}$, say a Sample Space is a set of ordered n-ples
\end{definition}

\begin{definition}
Let $S$ be a set. Say $\Omega \subseteq \mathcal{P}(S)$ is a $\sigma$-field whenever:
\begin{enumerate}
\item $\emptyset \in \Omega$
\item $\forall E \in \Omega, E^C \in \Omega$
\item If $\{E_n\}$ is a countable collection in $\Omega$, $\cup_n E_n \in \Omega$
\item If $\{E_n\}$ is a countable collection in $\Omega$, $\cap_n E_n \in \Omega$
\end{enumerate}
\end{definition}

\textbf{Ex-} Let $S=\{1, 2, 3, 4\}$. THen $\{\emptyset, \{1,2\}, \{3,4\}, S\}$ is a $\sigma$-field.

\begin{lemma}
Let $S$ be a set and $\Omega_1$, $\Omega_2$ be $\sigma$-fields on $S$. Then $\Omega_1 \cap \Omega_2$ is a $\sigma$-field on $S$
\end{lemma}
\begin{proof}
Clearly, $\emptyset \in \Omega_1 \cap \Omega_2$, as $\emptyset \in \Omega_1, \Omega_2$.
Next, let $E\in \Omega_1 \cap \Omega_2$.
THen $E \in \Omega_1 , \Omega_2$, so $E^C \in \Omega_1 , \Omega_2$.
Thus, $E^C \in \Omega_1 \cap \Omega_2$. \\
Finally, let $\{E_n \}$ be a countable collection of sets in $\Omega_1 \cap \Omega_2$. 
THen every $E_n$ is in both $\Omega_1$ and $\Omega_2$.
So $\cup_n E_n, \cap_n E_n \in \Omega_1, \Omega_2$.
Thus, $\cap_n E_n, \cup_n E_n \in \Omega_1 \cap \Omega_2$.
Therefore, $\Omega_1 \cap \Omega_2$ is a $\sigma$-field.
\end{proof}

\begin{definition}
Let $S$ be a set and $U \subseteq \mathcal{P}(S)$ be any collection of subsets of $S$.
Say the $\sigma$-field generated by $U$, $\mathscr{M}(U)$ is hte intersection of all $\sigma$-fields on $S$ which contain $U$.
\end{definition}

Note- $\mathscr{M} (U)$ is hte "smallest" $\sigma$-field that contains $U$. Inded, based on teh definition, if $U \subseteq \Omega$ and $\Omega$ is a $\sigma$-field then $\mathscr{M}(U) \subseteq \Omega$.

\begin{definition}
Let $\mathscr{T}$ be hte collection of all open sets in $\mathbb{R}$.
Say the Borel $\sigma$-field is $\mathcal{B} = \mathscr{M}(\mathscr{T})$.
\end{definition}

\begin{definition}
Let $S$ be a sample space and $\Omega$ a $\sigma$-field on $S$.
Say $E \in \Omega$ is an event.
\end{definition}

\begin{definition}
Let $S$ be a sample space with $\sigma$-field $\Omega$.
Say a function $p: \Omega \rightarrow \mathbb{R}$ is a probability function whenever
\begin{enumerate}
\item $\forall E \in \Omega, p(E) \geq 0$
\item $p(S)$ = 1
\item For any collection of pairwise disjoint events $\{ E_n \}$, $p(\cup_n E_n ) = \sum_n p(E_n )$
\end{enumerate}
\end{definition}

\begin{definition}
Let $S$ be a sample space, $\Omega$ a $\sigma$-field on $S$, and $p: \Omega \rightarrow \mathbb{R}$ a probability function.
Say $(S, \Omega, p)$ is a probability space. 
\end{definition}

\textbf{Ex-} Let $S = \{ HH, HT, TH, TT \}$, $\Omega = \mathcal{P}(S)$, and $p(E) = \frac{|E|}{|S|}$.
This creates a probability space, representing the outcome of flipping two fair coins.
In this case, $p(\{HH, HT \}) = \frac{|\{HH, HT \}|}{| \{ HH, HT, TH, TT \}|} = \frac{2}{4}=\frac{1}{2}$.

\begin{lemma}
Let $(S, \Omega, p)$ be a probability space and $E \in \Omega$.
Then $p(E^C) = 1 - p(E)$.
\end{lemma}
\begin{proof}
Notice $\{ E, E^C \}$ is a collection of pairwise disjoint sets in $\Omega$, so \\$p(E \cap E^C ) = p(E) + p(E^C )$.
Also, $E\cup E^C=S$, so $p(E) + p(E^C )= 1$.
\end{proof}
\begin{corrolary}
$p(\emptyset ) = 0$
\end{corrolary}

\begin{lemma}
A probability space is a measure space with total measure 1.
\end{lemma}
\begin{corrolary}
\textbf{Union Bound: }
$p(E_1 \cup E_2) \leq p(E_1) + p(E_2)$
\end{corrolary}

\begin{theorem}
\textbf{Principle of Inclusion-Exclusion:}
Let $(S, \Omega, p)$ be a probability space and $\{E_n\}$ be a collection of events.
Then
\begin{itemize}
\item $p(E_1 \cup E_2 ) = p(E_1) + p(E_2) - p(E_1 \cap E_2 )$
\item $p(E_1 \cup E_2 \cup E_3 ) = p(E_1) + p(E_2) + p(E_3)_ - p(E_1 \cap E_2 ) - p(E_1 \cap E_3) - \allowbreak p(E_2 \cap E_3 ) + p(E_1 \cap E_2 \cap E_3 )$
\item and so on, where we add $(-1)^{k-1}$ times the sum of hte probabilities of all $k$-sized intersections of the events
\end{itemize}
\end{theorem}
\begin{proof}
Let $E$ and $F$ be any two events and consider that $E=(E \cap F ) \cup (E \cap F^C )$ and $F = (E \cap F) \cup (E^C \cap F)$.
So $E \cup F = (E \cap F) \cup (E \cap F^C ) \cup (E^C \cap F)$.
Also, $\{E \cup F, E \cup F^C, E^C \cup F\}$ is a countable collection of disjoint sets.
So $p(E \cup F ) =  p(E \cap F) + p(E \cap F^C ) + p(E^C \cap F)$ and $p(E) + p(F) = 2 p(E \cap F) + p(E \cap F^C ) + p(E^C \cap F)$.
So $p(E) + P(F) = p(E\cup F) + p(E\cap F)$.
Now, we will show the PIE by induction. \\
In the base case, let $n=2$ be the number of events and let $E=E_1$ and $F=E_2$. Then we have established $p(E_1) + p(E_2) = p(E_1 \cap E_2 ) + p(E_1 \cup E_2)$.
Thus, this establishes the base case. \\
Now, suppose the PIE holds for some $n \in \mathbb{N}$.
We will let $E=\cup_{k=1}^n E_k$ and $F=E_{n+1}$.
Then $p(\cup_{k=1}^{n+1} E_k ) = p(\cup_{k=1}^n E_k ) + p(E_{n+1}) - p((\cup_{k=1}^n E_k) \cap E_{n+1}) =p(\cup_{k=1}^n E_k ) + p(E_{n+1}) - p(\cup_{k=1}^n (E_k \cap E_{n+1}))$.
By applying PIE to the first and third terms, we have the desired result.
Therefore, by induction, PIE holds for all $N$.
\end{proof}

\begin{definition}
Say a collection of events $\{ E_n \}$ is a cover when $S = \cup_n E_n$.
\end{definition}

\begin{definition}
Say a partition is a cover when all $E_n$ are pairwise disjoint.
\end{definition}

\begin{lemma}
If $\{E_n \} \searrow E$ then $p(E_n ) \searrow p(E)$.
\end{lemma}
\begin{proof}
Let $R_n = E_n \cap E_{n+1}^C$.
Then $\{R_k\}_{k=1}^\infty$ and $E$ together form a countable partition of $E_n$.
So $p(E_n) = p(E) + \sum_{k=n}^\infty p(R_k)$.
Because $p(E)$ and $p(E_n)$ are finite, the sum converges.
Let $\epsilon > 0$.
Then there is some $N \in \mathbb{N}$ such that for all $n \geq N$, $\sum_{k=n}^\infty p(R_k) \leq \epsilon$.
So $p(E) \leq p(E_n) \leq p(E) = \epsilon$.
\end{proof}
\begin{corrolary}
If $\{E_n\} \nearrow E$, $p(E_n) \nearrow p(E)$.
\end{corrolary}

\begin{theorem}
\textbf{Borel-Cantelli Lemma:}
Let $(S, \Omega, p)$ be a probability space, $\{ E_n \}$ a sequence of events, and $E = \limsup\limits_{n} E_n$.
If $\sum_{n=1}^\infty p(E_n) \leq \infty$ then $p(E) = 0$.
\end{theorem}
\begin{proof}
Recall $\limsup\limits_{n} E_n = \cap_{m \geq 1} \cup_{n \geq m} E_n$.
For all $m \geq 1$, $E \subseteq \cup_{n \geq m} p(E_n$.
Furthermore, this union induces a decreasing sequence of sets as $m$ increases.
So $p(E) \leq p(\cup_{n \geq m} E_n) \leq \sum_{n\geq m} p(E_n)$.
If $\sum_{n\geq 1} p(E_n) \leq \infty$ then $\sum_{n\geq m} p(E_n) \rightarrow 0$ as $m \rightarrow \infty$.
Thus, $p(E)=0$.
\end{proof}

\end{document}
